% generated by roam.py handbook export -- не редактировать, править заметки
% 20260918155412-выпуклость_множества_подуровня_выпуклои_функции.org (body 08743327dc0c)

\begin{theorem}[{Выпуклость множества подуровня выпуклой функции}]\label[theorem]{note:1CD538BF-6F40-472A-A670-D38F093E4DF1}
  Пусть $f: E \rightarrow \R$~--- \hyperref[note:1CA14F2C-637F-41DE-9E97-72075D19E97F]{выпуклая функция}, $c \in \R$, $E$~--- \hyperref[note:316B5BFA-8B9F-4267-9086-D6593AAB5F9B]{выпуклое}
  Тогда $\mathrm{Lev}_f(c) = \left\{ x \in E \  |\  f(x) \le c  \right\}$~--- \hyperref[note:316B5BFA-8B9F-4267-9086-D6593AAB5F9B]{выпуклое множество}
\end{theorem}

\begin{proof}
  Будем проверять определение выпуклого множества.

  Возьмем произвольные $x, y \in \mathrm{Lev}_f(c), \lambda \in [0;1]$.
  Заметим, что $\mathrm{Lev}_f(c) \subset E \Rightarrow x,y \in E$
  Знаем, что $\lambda x + (1-\lambda) y \in E$ \hyperref[note:316B5BFA-8B9F-4267-9086-D6593AAB5F9B]{(по выпуклости)}

  По выпуклости $f$: $f(\lambda x + (1-\lambda)y) \le \lambda f(x) + (1-\lambda)f(y)$

  \[
  x,y \in \mathrm{Lev}_f(c) \Rightarrow \lambda f(x) + (1-\lambda)f(y) \le \lambda c + (1-\lambda)c = c
  \]

  Но тогда, $f(\lambda x + (1-\lambda)y) \le c \Rightarrow (\lambda x + (1-\lambda)y) \in \mathrm{Lev}_{f}(c)$
  Итого $\mathrm{Lev}_f(c)$~--- выпуклое \hyperref[note:316B5BFA-8B9F-4267-9086-D6593AAB5F9B]{по определению}.
\end{proof}
